%==============================================================================
%  CAPÍTULO — LAS GARANTÍAS REALES EN EL CONCURSO
%==============================================================================
\chapter{Las garantías reales en el concurso}
\label{cap:garantias-reales}

\fichaCapitulo{El acreedor garantizado frente al procedimiento colectivo.}{Segunda y tercera clase
de créditos, ejecución individual y concursal, el derecho de opción del acreedor prendario y
retencionario, la purga de la hipoteca y la posición del garante frente al \discharge{}.}

%==============================================================================
\bloqueTeorico
%==============================================================================

\subsection{Por qué la garantía real resiste el concurso}

El acreedor garantizado ocupa en el concurso una posición singular que conviene fundar antes de
describirla. La prenda y la hipoteca no confieren un mero privilegio de pago: confieren un
\destacar{derecho real} sobre una cosa determinada, oponible \emph{erga omnes} y anterior a la
insolvencia. Ese derecho no nace del concurso ni se subordina a él; el procedimiento colectivo lo
encuentra ya constituido y debe respetarlo.

La razón es a la vez dogmática y económica. Dogmática, porque el derecho real sigue a la cosa con
independencia de quién sea su dueño o de su estado patrimonial. Económica, porque ---como se expuso
al tratar la teoría del acuerdo hipotético (capítulo \ref{cap:teoria-economica})--- si la garantía
pudiera desconocerse en el concurso, precisamente cuando importa, dejaría de cumplir su función y el
crédito garantizado se encarecería para todos los deudores. Respetar la garantía en la insolvencia
es la condición de su existencia fuera de ella.

\begin{foco}[La garantía no es un privilegio más]
Conviene no confundir la preferencia de la garantía real con los privilegios de primera clase del
\artcc{2472}. El crédito laboral o el fiscal preferente se pagan \emph{antes} con el producto
general del activo, pero no tienen una cosa afectada: su preferencia es de rango. El acreedor
hipotecario o prendario, en cambio, tiene \destacar{una cosa determinada} sujeta a su derecho, y se
paga con ella con preferencia a casi todos. La diferencia es la que media entre un mejor lugar en la
fila y tener asignado un bien concreto.
\end{foco}

\subsection{La ubicación en la prelación: segunda y tercera clase}

El \artcc{2472} agota la primera clase; las garantías reales viven en las clases siguientes:

\begin{description}[leftmargin=0pt,style=nextline,font=\sffamily\bfseries]
  \item[Segunda clase (\artcc{2474})] Los créditos del posadero y del acarreador sobre los efectos
  que tienen en su poder, y ---la hipótesis relevante--- el \destacar{acreedor prendario sobre la
  prenda}. Son preferencias \emph{especiales}: recaen sobre bienes muebles determinados.
  \item[Tercera clase (\artcc{2477})] Los créditos \destacar{hipotecarios}, que se pagan con el
  producto de la finca hipotecada según el orden de sus inscripciones.
\end{description}

La relación entre estas clases y la primera se rige por el \artcc{2476}: sobre una misma especie,
los créditos de segunda clase \emph{excluyen} a los de primera; pero si los demás bienes del deudor
son insuficientes para cubrir la primera clase, ésta se paga con la especie afectada en cuanto al
déficit. Es la regla que permite a los créditos laborales y previsionales alcanzar, en subsidio, el
producto de los bienes gravados.

%==============================================================================
\bloqueNormativo
%==============================================================================

\subsection{El derecho de opción del acreedor de segunda clase}

El \art{242} confiere al acreedor prendario y al retencionario ---cuyo derecho de retención haya
sido judicialmente declarado--- un \destacar{derecho de opción} de gran importancia práctica: pueden
optar por

\begin{enumerate}
  \item \textbf{Ejecutar individualmente} el bien gravado, iniciando o continuando el procedimiento
  ante el tribunal que conoce de la liquidación ---previa acumulación--- y asegurando siempre los
  créditos de mejor derecho; o
  \item \textbf{Concurrir al procedimiento concursal} verificando su crédito y su preferencia, y
  cobrando con el producto de la realización que administre el liquidador.
\end{enumerate}

La ley reconoce, además, una facultad al liquidador: \destacar{exigir la entrega de la cosa} dada en
prenda o retenida, siempre que pague la deuda garantizada o deposite a la orden del tribunal su valor
estimativo, sobre el cual se hará efectiva la preferencia. Es la vía para que la masa recupere un
bien cuya realización conjunta con otros ---como unidad económica--- rendiría más que su ejecución
aislada.

\begin{practica}[Cuándo conviene ejecutar y cuándo concurrir]
La elección no es neutra. Ejecutar individualmente da control sobre los tiempos y la forma de
realización, pero obliga a asegurar los créditos de mejor derecho y expone al acreedor a las cargas
del procedimiento ejecutivo. Concurrir al concurso descarga esa gestión en el liquidador, pero
somete el bien a la estrategia de realización de la masa ---que puede diferir la venta para
integrarlo en una unidad económica---. La regla práctica: cuando el bien tiene un mercado líquido y
propio, ejecutar; cuando su valor depende de venderse con otros activos, concurrir y negociar.
\end{practica}

\subsection{Los acreedores hipotecarios y el concurso especial}

El \art{243} remite el pago de los acreedores hipotecarios a los \artcc{2477} a 2480. De ese
régimen derivan las reglas que estructuran su posición:

\paragraph{El orden de las inscripciones.} Los hipotecarios se pagan con el producto de la finca en
el orden de sus inscripciones (\artcc{2477}), de modo que la primera hipoteca excluye a la segunda y
así sucesivamente. La fecha de la inscripción ---no la del crédito--- fija el rango.

\paragraph{El concurso especial de acreedores hipotecarios (\artcc{2477}).} Los hipotecarios de una
misma finca pueden provocar su realización \destacar{sin aguardar el resultado del concurso
general}, formando un concurso especial en que se les paga de inmediato según su rango, con la carga
de rendir caución por lo que pueda corresponder a los créditos de primera clase en el déficit del
resto del patrimonio. Es la manifestación más nítida de la fuerza de la garantía: el hipotecario no
espera.

\paragraph{La posposición de la primera clase (\artcc{2478}).} Si los demás bienes son insuficientes
para cubrir los créditos de primera clase, éstos se pagan con el valor de la finca hipotecada, con
preferencia a los hipotecarios ---que, habiendo cauccionado o pagado, se subrogan---. La superpreferencia
laboral y fiscal alcanza, así, incluso al bien hipotecado, pero sólo en subsidio.

\begin{alerta}[La hipoteca no se extingue por la sola apertura del concurso]
Un error frecuente es suponer que la liquidación «purga» automáticamente las hipotecas. No es así:
la garantía subsiste y sigue a la finca. Su extinción requiere el pago del crédito garantizado o la
\destacar{purga de la hipoteca} por la realización en pública subasta con citación de los acreedores
hipotecarios (\artcc{2428} y \artcpc{492}). El adquirente que compra en la realización concursal sin
que se hayan purgado las hipotecas las recibe \emph{con} el gravamen.
\end{alerta}

\subsection{La purga de la hipoteca en la realización concursal}

La \destacar{purga} es el mecanismo por el cual la venta forzada de la finca extingue las hipotecas
que la gravan, permitiendo transferirla libre al adquirente. Sus condiciones, tomadas del \artcc{2428}
y del \artcpc{492}:

\begin{enumerate}
  \item Que la finca se realice en \destacar{pública subasta} ordenada por el tribunal.
  \item Que los acreedores hipotecarios sean \destacar{citados} personalmente.
  \item Que el producto se consigne a disposición del tribunal para pagar a los hipotecarios según
  su rango.
\end{enumerate}

Purgada la hipoteca, el rango de los acreedores no pagados se traslada al precio: cobran sobre el
producto en el mismo orden en que habrían cobrado sobre la finca. El acreedor hipotecario de grado
posterior citado y no pagado por insuficiencia del precio \emph{pierde} su hipoteca; el no citado la
conserva.

\begin{practica}[La citación de los hipotecarios es el punto crítico]
Para el liquidador que realiza un inmueble gravado, la citación válida de \emph{todos} los
acreedores hipotecarios es la actuación que blinda la venta: sin ella no hay purga, y el inmueble se
vende con las hipotecas, lo que deprime su precio o frustra la venta. Para el acreedor hipotecario,
a la inversa, verificar que fue citado ---y comparecer--- es la condición de conservar su rango sobre
el precio. Una citación defectuosa es fuente de nulidades y de responsabilidad del liquidador.
\end{practica}

%==============================================================================
\bloquePractico
%==============================================================================

\subsection{El garante real y personal frente al \discharge{}}

La posición de quien garantizó la deuda del concursado ---fiador, codeudor solidario, avalista o
tercero que hipotecó o prendó un bien propio--- merece tratamiento propio, porque concentra uno de
los errores de asesoría más costosos.

Como se expuso al tratar la extinción de saldos (capítulo \ref{cap:liquidacion-ordinaria}), el
\art{255} es categórico: la descarga del deudor \destacar{no afecta los derechos del acreedor contra
los garantes}. El fiador, el avalista y el tercero constituyente de garantía real siguen obligados
por el total, \emph{no} pueden invocar la extinción que benefició al deudor principal y ---dato
decisivo--- \emph{no} pueden subrogarse contra el deudor descargado ni exigirle reembolso por lo que
paguen.

\begin{alerta}[El aval sobrevive a la descarga del avalado]
La deuda que la ley extinguió para el deudor sigue \emph{viva} para quien la garantizó. Quien avaló
a una empresa o a un familiar que luego se liquidó responde por el total, y si paga no puede repetir
contra el rehabilitado. Es una advertencia que no puede faltar en la asesoría ---tanto al deudor,
que debe saber que no libera a sus garantes, como al garante, que debe saber que su obligación no se
descarga con la del principal---.
\end{alerta}

\subsection{Cuadro de posiciones del acreedor garantizado}

\begin{table}[htbp]
\centering
\small
\caption{El acreedor con garantía real en el procedimiento concursal}
\label{tab:cap11b-garantias}
\begin{tabularx}{\textwidth}{@{}>{\raggedright\arraybackslash}p{0.24\textwidth}
                                 >{\raggedright\arraybackslash}p{0.24\textwidth} X@{}}
\toprule
\textbf{Cuestión} & \textbf{Prendario / retencionario} & \textbf{Hipotecario} \\
\midrule
Clase (\artcc{}) & Segunda (2474) & Tercera (2477) \\
\addlinespace
Ejecución individual & Sí, con acumulación (\art{242}) & Sí, concurso especial (\artcc{2477}) \\
\addlinespace
Rango & Sobre la cosa mueble & Orden de las inscripciones \\
\addlinespace
Cede ante primera clase & En subsidio, si falta activo (\artcc{2476}) & En subsidio (\artcc{2478}) \\
\addlinespace
El liquidador puede & Exigir entrega pagando o depositando el valor (\art{242}) & Realizar con
purga previa citación \\
\addlinespace
Extinción de la garantía & Pago o realización & Pago o purga (\artcc{2428}, \artcpc{492}) \\
\bottomrule
\end{tabularx}
\end{table}

\begin{foco}[La garantía real es la mejor posición, pero no es inmune]
El acreedor con garantía real es el mejor situado del concurso: tiene una cosa afectada, puede
ejecutarla por separado y su rango es anterior. Pero no es intocable ---cede en subsidio ante la
superpreferencia laboral y fiscal, su garantía puede revocarse si se constituyó en el período
sospechoso por una deuda previa (capítulo \ref{cap:revocatorias}), y su hipoteca se purga en la
realización si es citado---. Conocer con precisión esos límites es lo que separa al acreedor
garantizado que cobra del que descubre, tarde, que su garantía valía menos de lo que creía.
\end{foco}
